\chapter{Durchführung}
\section{Caption Tests}
\begin{figure}[htbp!]
    \centering
    \includegraphics{img/dhbwlogo.png.png}
    \caption{Here is a really dragged out capiton. I mean it is really streched thin. For how much longer could this caption go on. I'll hazard a guess and say not much longer.}
    \label{fig1}
\end{figure}

\captionsetup[figure]{format=hang, justification=raggedright, labelsep=colon, labelfont=normal}

\begin{figure}[htbp!]
    \centering
    \includegraphics{img/dhbwlogo.png.png}
    \caption{Here is a really dragged out capiton. I mean it is really streched thin. For how much longer could this caption go on. I'll hazard a guess and say not much longer.}
    \label{fig2}
\end{figure}

\captionsetup[figure]{justification=justified}

\begin{figure}[htbp!]
    \centering
    \includegraphics{img/dhbwlogo.png.png}
    \caption{Here is a really dragged out capiton. I mean it is really streched thin. For how much longer could this caption go on. I'll hazard a guess and say not much longer.}
    \label{fig3}
\end{figure}

\captionsetup[figure]{justification=centering}

\begin{figure}[htbp!]
    \centering
    \includegraphics{img/dhbwlogo.png.png}
    \caption{Here is a really dragged out capiton. I mean it is really streched thin. For how much longer could this caption go on. I'll hazard a guess and say not much longer.}
    \label{fig4}
\end{figure}

\captionsetup[figure]{justification=centerlast, format=plain}

\begin{figure}[htbp!]
    \centering
    \includegraphics{img/dhbwlogo.png.png}
    \caption{Here is a really dragged out capiton. I mean it is really streched thin. For how much longer could this caption go on. I'll hazard a guess and say not much longer.}
    \label{fig5}
\end{figure}

\captionsetup[figure]{justification=justified, format=plain}

\begin{figure}[htbp!]
    \centering
    \includegraphics{img/dhbwlogo.png.png}
    \caption{Here is a really dragged out capiton. I mean it is really streched thin. For how much longer could this caption go on. I'll hazard a guess and say not much longer.}
    \label{fig6}
\end{figure}

\captionsetup[figure]{justification=justified, format=plain, labelsep=quad}

\begin{figure}[htbp!]
    \centering
    \includegraphics{img/dhbwlogo.png.png}
    \caption{Here is a really dragged out capiton. I mean it is really streched thin. For how much longer could this caption go on. I'll hazard a guess and say not much longer.}
    \label{fig7}
\end{figure}

\captionsetup[figure]{format=hang, justification=justified, format=plain, labelsep=quad, labelfont = bf}

\begin{figure}[htbp!]
    \centering
    \includegraphics{img/dhbwlogo.png.png}
    \caption{Here is a really dragged out capiton. I mean it is really streched thin. For how much longer could this caption go on. I'll hazard a guess and say not much longer.}
    \label{fig8}
\end{figure}

\cleardoublepage

\begin{figure}[htbp!]
    \centering
    \includegraphics[width=0.5\textwidth]{img/dhbwlogo.png}
    \caption{Short caption}
    \label{fig9}
\end{figure}

Here a table

\begin{table}[!htbp]
    \centering
    \begin{tabular}{@{}llll@{}}
    \toprule
    \multicolumn{2}{c}{elektrische Größe} & \multicolumn{2}{c}{thermische Größe} \\ \cmidrule(r){1-2} \cmidrule(l){3-4} 
    Größe & Einheit & Größe & Einheit \\ \midrule
    Potenzial $\varphi$ & $\unit{\V}$ & Temperatur $T$ & \unit{\K}\\
    Strom $I$ & \unit{\A} & Wärmestrom $\dot Q$ & \unit{\W} \\
    Widerstand $R$ & $\unit{\ohm} = \unit{\V\per\A}$ & thermischer Widerstand $R_{th}$ & \unit{\K\per\W} \\ 
    Ladung $Q$ & \unit{\C} & Wärmemenge $Q$ & \unit{\J} \\
    Leitfähigkeit $\kappa$ & \unit{\per\ohm\per\meter} & thermische Leitfähigkeit $\lambda$ & \unit{\W\per\m\per\K} \\ \bottomrule
    \end{tabular}
    \caption{Analogien von thermischen und elektrischen Größen}
    \label{tab:electricalThermalAnalogy}
    Quelle: In Anlehnung an \cite[S. ~226]{Tietze2010}
\end{table}


\begin{equation}
    x_1^2+x_2^2+x_3^2 = |\vec{x}|
\end{equation}
mit
\begin{symboltable}
    $\vec{x}$:  & x-Vektor\\
    $x_1$:      & x-Komponente von $\vec{x}$\\
    $x_2$:      & y-Komponente von $\vec{x}$\\
    $x_2$:      & z-Komponente von $\vec{x}$\\
\end{symboltable}
Hier gehts weiter mit Text.